\documentclass[12pt, a4paper]{article}

% --- Encoding & Fonts ---
\usepackage[utf8]{inputenc}
\usepackage[T1]{fontenc}
\usepackage{lmodern}

% --- Mathematics ---
\usepackage{amsmath}
\usepackage{amssymb}
\usepackage{amsthm}
\usepackage{mathtools}

% --- Layout ---
\usepackage[margin=1in]{geometry}

% --- Lists ---
\usepackage{enumitem}

% --- Tables ---
\usepackage{booktabs}
\usepackage{array}

% --- Graphics ---
\usepackage{graphicx}

% --- Headers / Footers ---
\usepackage{fancyhdr}
\pagestyle{fancy}
\fancyhf{}
\fancyhead[L]{\leftmark}
\fancyhead[R]{\thepage}
\renewcommand{\headrulewidth}{0.4pt}

% --- Hyperlinks & Cross-references ---
\usepackage[colorlinks=true, linkcolor=blue, urlcolor=blue, citecolor=blue]{hyperref}

% --- Misc ---
\usepackage{xcolor}

% --- Custom commands (safety definitions) ---
\providecommand{\operatorname}[1]{\mathop{\mathrm{#1}}\nolimits}

\title{Learning Module: Sdos2Projectmanagement}
\author{}
\date{\today}

\begin{document}

\maketitle
\thispagestyle{empty}
\newpage

\tableofcontents
\newpage

% ============================================================
% SECTION  Project Management Overview
% ============================================================

% !!!  REVIEW FLAG: Project Management Overview  !!!
% [undefined_term] Define or link to an explanation of the burndown chart in a previous module or include a brief definition in this section.
% [undefined_term] Introduce earned value analysis and explain how it relates to project monitoring earlier in the learning path or provide a concise definition here.
% !!! END REVIEW FLAGS !!!
\section{Project Management Overview}

\subsection{Before We Begin}
Software development typically follows a sequence of stages: first gathering requirements, then designing the system, followed by coding the solution, and finally testing it.  Each stage produces work that must be coordinated with the others.  Management activities such as planning, tracking progress, and handling risks are layered on top of these stages to keep the effort focused, on schedule, and within budget.  Understanding this overlay helps us see where projectlevel planning and monitoring fit into the overall development effort.  

\subsection{Motivation}
Planning a roadtrip: you decide on destinations, map the route, and adjust the itinerary as traffic or weather changes.  In software projects, a similar process is needed to decide what to build, how to build it, and how to respond when conditions shift.  Project management provides the structure that keeps a development effort organized, on track, and resilient to surprises.  

\subsection{Intuition}
Think of a software project as a journey.  The \textbf{Planning} phase is like deciding where you want to go, what stops you will make, and estimating how long each leg will take.  Once the journey begins, the \textbf{Monitoring and Control} phase is the continual checking of your position, speed, and resources, allowing you to steer back on course when you encounter delays or obstacles.  The plan is not a static map; it is a living guide that is updated as new information appears, such as traffic jams (risks) or detours (change requests).  

\subsection{Formal Definition}
\textbf{Project Management} is the organized effort to guide a software development effort from start to finish.  It consists of two primary phases:

\begin{enumerate}
    \item \textbf{Planning}: This initial activity creates a \emph{Project Management Plan}.  The plan contains:
    \begin{itemize}
        \item \textit{Cost and schedule estimation} --- predicts the monetary resources and timeline needed.
        \item \textit{Quality planning} --- defines the standards and activities that ensure the product meets requirements.
        \item \textit{Monitoring plans} --- describes how progress will be measured and reported.
        \item \textit{Risk management plans} --- identifies potential problems and outlines responses.
    \end{itemize}
    The plan is deliberately flexible; it serves as a baseline that can be revised as the project proceeds.
    
    \item \textbf{Monitoring and Control}: This phase occupies the entire lifespan of the project.  It continuously compares actual performance against the baseline, adjusts assignments, revises estimates, and updates the \emph{Project Management (Execution) Plan}---the detailed, actionable document that states who does what, by when, and how quality will be assured.  Because realworld conditions evolve, the Execution Plan must be \emph{dynamic}, meaning it is regularly refreshed to reflect new information, change requests, or emerging risks.
\end{enumerate}

\subsection{Worked Example}
\textbf{Problem statement:}  Create a onepage Project Management Plan for a small mobile todolist app and show how the plan is updated after a change request that adds a ``calendar integration'' feature.

\textbf{Step 1 -- Identify core components of the initial plan.}

\begin{tabular}{|p{3cm}|p{9cm}|}
\hline
\textbf{Component} & \textbf{Initial Content} \\
\hline
Cost Estimate & Development: \$8,000; Design: \$2,000; Testing: \$2,000; Total: \$12,000. \\
\hline
Schedule & Weeks 1--2: Requirements; Weeks 3--5: Design; Weeks 6--10: Coding; Weeks 11--12: Testing; Week 13: Release. \\
\hline
Quality Activities & Code reviews each sprint; Automated unit tests covering 80\% of code; Useracceptance testing (UAT) in week 12. \\
\hline
Monitoring Approach & Weekly status meetings; Burndown chart for effort remaining; Earned value analysis biweekly. \\
\hline
Top Three Risks & R1: Underestimated UI design effort; R2: Thirdparty library compatibility; R3: Delayed feedback from target users. \\
\hline
\end{tabular}

\textbf{Step 2 -- Receive a change request.}  The client asks to add ``calendar integration'' so users can view tasks in their device calendar.

\textbf{Step 3 -- Assess impact.}
\begin{itemize}
    \item \textbf{Cost} -- Additional development effort estimated at \$3,000; testing at \$1,000.
    \item \textbf{Schedule} -- Adds two weeks of coding (weeks 11--12) and pushes release to week 15.
    \item \textbf{Quality} -- New integration tests needed; update UAT criteria.
    \item \textbf{Risks} -- New risk R4: API changes in the calendar service could cause delays.
\end{itemize}

\textbf{Step 4 -- Update the Project Management Plan.}

\begin{tabular}{|p{3cm}|p{9cm}|}
\hline
\textbf{Component} & \textbf{Revised Content} \\
\hline
Cost Estimate & Original \$12,000 + \$4,000 (new feature) = \$16,000. \\
\hline
Schedule & Weeks 1--2: Requirements; Weeks 3--5: Design; Weeks 6--10: Core coding; Weeks 11--12: Calendar integration coding; Weeks 13--14: Testing (including integration); Week 15: Release. \\
\hline
Quality Activities & Add integration test suite for calendar API; extend UAT to cover calendar sync scenario. \\
\hline
Monitoring Approach & Include a milestone ``Calendar Integration Complete'' in the burndown chart; track risk R4 in weekly risk log. \\
\hline
Top Three Risks & Update list: R1, R2, R3, plus R4 (new). \\
\hline
\end{tabular}

\textbf{Final answer:}  The revised onepage plan now reflects the higher cost, extended schedule, additional quality checks, and a new risk, demonstrating that the Execution Plan evolves in response to change requests.

\textbf{What this shows us:}  Project management is not a onetime paperwork task; the plan must be revisited and adjusted whenever the project's reality changes, ensuring the team stays aligned with goals, resources, and risk exposure.

\subsection{Key Takeaway}
Project management supplies a flexible plan and ongoing monitoring that keep a software project coordinated, on schedule, and responsive to new information.

\subsection*{Questions}

% --- RECALL (1) ---
\begin{enumerate}
\item \textbf{Recall:} List the two primary phases of project management that are described in the section above.

\newpage

% --- UNDERSTAND (2) ---
\item \textbf{Understand:} In your own words, explain why the Project Management Execution Plan is described as ``dynamic.'' Be sure to refer to how the plan is used during the Monitoring and Control phase.

\newpage

% --- UNDERSTAND (3) ---
\item \textbf{Understand:} Describe how a dynamic Execution Plan helps a project team respond to a change request (e.g., a new feature or a newly identified risk). Use examples from the worked example in the theory.

\newpage

% --- APPLY (2) ---
\item \textbf{Apply:} A small elearning mobile app is under development. Midproject the client requests a ``pushnotification reminder'' feature.  
Identify which components of the Project Management Plan (Cost Estimate, Schedule, Quality Activities, Monitoring Approach, and Top Risks) must be updated to accommodate this change. List each component that needs revision and give a brief reason for the update.

\newpage

% --- APPLY (2) ---
\item \textbf{Apply:} Using the information from the worked example (the todolist app), create a short revised ``Cost Estimate'' and ``Schedule'' table that reflects the addition of the pushnotification reminder feature. Assume the new feature adds \$2,500 in development cost, \$800 in testing cost, and requires two additional weeks of coding after week 10. Show your revised numbers in the same tabular format as the example.

\newpage

% --- ANALYSE (1) ---
\item \textbf{Analyse:} Compare two project plans for the same mobile app:  
1) a plan that treats the Project Management Plan as a static document (no changes after the initial baseline), and  
2) a plan that treats the Project Management Plan as dynamic (updated whenever a change request occurs).  
Discuss the likely impact of each approach on project cost, schedule, quality, and overall success, drawing on the concepts presented in the section.

\newpage

% --- EVALUATE (1) ---
\item \textbf{Evaluate:} Below is a brief Project Management Plan for a hypothetical ecommerce website (excerpt):

\begin{tabular}{|p{3cm}|p{9cm}|}
\hline
\textbf{Component} & \textbf{Content} \\
\hline
Cost Estimate & Development: \$15,000; Design: \$5,000; Testing: \$4,000; Total: \$24,000. \\
\hline
Schedule & Weeks 1--3: Requirements; Weeks 4--6: Design; Weeks 7--12: Development; Weeks 13--15: Testing; Week 16: Launch. \\
\hline
Quality Activities & Code reviews at end of each sprint; Manual UI testing before launch. \\
\hline
Monitoring Approach & Biweekly status reports; No burndown chart or earnedvalue analysis. \\
\hline
Top Three Risks & R1: Payment gateway integration; R2: High traffic spikes; R3: Content delivery delays. \\
\hline
\end{tabular}

Critique this plan by noting at least three strengths and three weaknesses, focusing on how well it addresses planning, ongoing monitoring, and risk management. Suggest concrete improvements that would make the plan more dynamic and effective.

\end{enumerate}

\newpage
% ============================================================
% SECTION  Effort Estimation
% ============================================================

\subsection*{Notation Used in This Section}
\begin{tabular}{@{}ll@{}}
\toprule
Symbol & Meaning \\ \midrule
    $Effort$ & Total effort in personmonths \\
    $size$ & Software size measured in KLOC \\
    $a$ & Effort scaling constant derived from historical data \\
    $b$ & Exponent reflecting economies or diseconomies of scale \\
    $E$ & Effort (personmonths) used in schedule formula \\
    $S$ & Schedule duration in months \\
\bottomrule
\end{tabular}
\vspace{1em}

% !!!  REVIEW FLAG: Effort Estimation  !!!
% [undefined_term] Introduce a brief definition of the COCOMO model early in the formal definition section, explaining that it is a parametric effort estimation model with cost drivers.
% [undefined_term] Define this cost driver explicitly, e.g., a multiplier that reflects the effort increase for meeting higher reliability requirements.
% [undefined_term] Define this cost driver as the multiplier representing the skill level of the development team and its impact on productivity.
% [undefined_term] Define this cost driver as the multiplier reflecting the effectiveness of software tools available for development and maintenance.
% [knowledge_assumption] Add a concise description of the Cone of Uncertainty, showing how it depicts decreasing estimate variability with increased information.
% !!! END REVIEW FLAGS !!!
\section{Effort Estimation}

\subsection{Before We Begin}
KLOC stands for ``thousands of lines of code'' and is used as a simple way to describe how large a software system is.  
A personmonth is a unit that represents the amount of work one person can accomplish in one month; it is the standard measure for effort in software projects.  
Both size (KLOC) and effort (personmonths) are needed to turn a vague idea of a project into a concrete plan.  
Understanding these two ideas lets us plug numbers into the effortestimation formulas that follow.

\subsection{Motivation}
Estimating how many gallons of paint you need for a house before you know the exact wall dimensions; you start with a rough guess and refine it as you measure each room.  
Project managers need a way to predict how much work a software project will require so they can set realistic schedules and budgets.  
Early guesses are often far off because many details are still unknown, so a systematic estimation approach helps reduce surprise later on.

\subsection{Intuition}
Think of effort estimation as estimating the amount of ``work fuel'' a project will consume.  
First, we look at the size of the software---how many lines of code we expect to write.  
Then we adjust that raw size by factors that capture how hard the work will be, such as how reliable the system must be or how experienced the team is.  
The resulting number tells us roughly how many personmonths of labor the project will need, guiding planning and resource allocation.

\subsection{Formal Definition}
Let \textbf{Effort} denote the total amount of work required, measured in personmonths.  
Let \textbf{size} denote the software size measured in KLOC (thousands of lines of code).  
Let \textbf{a} be a scaling constant that reflects the typical effort observed in past projects of similar type.  
Let \textbf{b} be an exponent that captures how effort grows with size; values of \(b\) greater than 1 indicate that larger projects become disproportionately harder, while values less than 1 suggest economies of scale.

The basic topdown effort estimate is  

\[
\text{Effort} = a \times \text{size}^{\,b}.
\]

In the COCOMO model this raw estimate is further refined by multiplying with several \emph{cost drivers} (e.g., required reliability, analyst capability, tool support).  The final adjusted effort is the product of the raw estimate and all relevant driver multipliers.  Note that the effect of size is moderated by the exponent \(b\); a larger code base does not always mean a proportionally larger effort, and the driver multipliers do not all have the same influence.

\subsection{Worked Example}
\textbf{Problem statement:}  
A project is estimated to have \(\text{size}=120\) KLOC.  The model constants are \(a=2.5\) and \(b=1.12\).  Three illustrative costdriver ratings are: required reliability \(=1.15\), analyst capability \(=0.90\), and tool support \(=1.00\).  Compute the raw effort, apply the driver multipliers, and obtain the final adjusted effort.  

\textbf{Step 1: Compute the raw effort}  

\[
\text{Effort}_{\text{raw}} = a \times \text{size}^{\,b}
= 2.5 \times 120^{1.12}.
\]

Calculate \(120^{1.12}\):
\[
\ln 120 \approx 4.7875,\qquad
1.12 \times \ln 120 \approx 5.363,
\]
\[
120^{1.12}=e^{5.363}\approx 213.6.
\]

Thus
\[
\text{Effort}_{\text{raw}} = 2.5 \times 213.6 \approx 534.0\ \text{personmonths}.
\]

\textbf{Step 2: Apply the first driver (required reliability)}  

\[
\text{Effort}_{1}= 534.0 \times 1.15 = 614.1\ \text{personmonths}.
\]

\textbf{Step 3: Apply the second driver (analyst capability)}  

\[
\text{Effort}_{2}= 614.1 \times 0.90 = 552.69\ \text{personmonths}.
\]

\textbf{Step 4: Apply the third driver (tool support)}  

\[
\text{Effort}_{\text{final}} = 552.69 \times 1.00 = 552.69\ \text{personmonths}.
\]

\textbf{Final answer:} The adjusted effort estimate is approximately 553 personmonths (rounded to the nearest whole month).  

\textbf{What this shows us:} Even with a modest size, the choice of costdriver values can increase or decrease the effort substantially, illustrating why both the size formula and the driver multipliers are essential for realistic planning.

\subsection{Key Takeaway}
Effort estimation combines a sizebased formula (\(\text{Effort}=a\times\text{size}^b\)) with adjustment factors (cost drivers) to produce a realistic prediction of the personmonths needed for a software project.

\subsection*{Questions}

\textbf{Recall Question:}  
What unit is commonly used to express effort estimates in software projects?  

\newpage

\textbf{Understanding Question 1:}  
Explain in your own words why the range of possible effort estimates narrows as a project progresses from the feasibility phase to the detailed design phase.  

\newpage

\textbf{Understanding Question 2:}  
Describe how the inclusion of costdriver multipliers in the COCOMO model refines a raw effort estimate that is based solely on size. What role do these drivers play in improving estimation accuracy?  

\newpage

\textbf{Application Question 1 (Numerical):}  
Using the topdown effortestimation formula \(\text{Effort}=a \times \text{size}^{\,b}\), compute the effort for a project with the following parameters:  
\[
\text{size}=80\ \text{KLOC},\qquad a=2.0,\qquad b=1.05.
\]  
Show all steps of your calculation and give the final answer in personmonths (rounded to one decimal place).  

\newpage

\textbf{Application Question 2:}  
A project has a size of 80 KLOC and the same constants \(a=2.0\) and \(b=1.05\) as in the previous question.  
Assume three costdriver multipliers: required reliability = 1.10, analyst capability = 0.95, and tool support = 1.00.  
First compute the raw effort using the formula, then apply the three multipliers sequentially to obtain the final adjusted effort. Present each intermediate result.  

\newpage

\textbf{Analysis Question:}  
Consider the same project (size = 80 KLOC, \(a=2.0\), \(b=1.05\)).  
Two alternative sets of costdriver multipliers are proposed:

\begin{itemize}
  \item \textbf{Set A}: reliability = 1.20, analyst capability = 0.85, tool support = 1.10
  \item \textbf{Set B}: reliability = 1.05, analyst capability = 0.95, tool support = 1.00
\end{itemize}

Calculate the final adjusted effort for each set and compare the results. Which set of drivers would likely lead to a more realistic project schedule, and why? Discuss the impact of higher reliability demands versus higher analyst capability on the estimate.  

\newpage

\textbf{Evaluation Question:}  
A project manager decides to use a bottomup estimation approach for a brandnew product that has no historical data or similar past projects. Critically evaluate this decision. Identify the potential risks and shortcomings of applying a bottomup method in this context, and propose at least two ways the estimation process could be strengthened (e.g., incorporating expert judgment, using analogous projects, applying parametric models, etc.).  

\newpage

\newpage
% ============================================================
% SECTION  Scheduling
% ============================================================

\subsection*{Notation Used in This Section}
\begin{tabular}{@{}ll@{}}
\toprule
Symbol & Meaning \\ \midrule
    $S$ & Schedule duration in months \\
    $E$ & Effort in personmonths \\
\bottomrule
\end{tabular}
\vspace{1em}

\section{Scheduling}
\subsection{Before We Begin}
In the previous section we learned how to estimate the amount of work a project requires, expressed as effort in person--months.  An effort estimate tells us how many person--months of labor are needed to complete all the work.  Because a calendar schedule tells us how long the work will stretch on the clock, we must translate effort into time.  The schedule formula uses the effort value as its input, so knowing the effort is essential before we can compute a schedule.  With this connection in mind we can now look at how overall and detailed schedules are built.

\section{Motivation}
Planning a wedding timeline: you first set the wedding date (overall schedule) and then schedule every individual activity like cake tasting, dress fitting, and invitations (detailed schedule).  An overall schedule gives stakeholders a clear date for project completion and shows how the total effort maps to calendar months.  A detailed schedule then breaks that period into concrete tasks so the team knows exactly what to do each day.  Together they solve the problem of turning abstract effort numbers into a practical, daybyday plan.

\subsection{Intuition}
Think of the overall schedule as the silhouette of a project: it outlines the major milestones and the final delivery date, based mainly on how much work must be done.  The detailed schedule fills that silhouette with the individual brush strokes---short tasks that a person can finish in a few days---assigning each to a specific person and time slot.  Stretching the silhouette (adding time) is easy; squeezing it (doing the same work faster) usually requires more resources and costs more.

\subsection{Formal Definition}
Let \(S\) denote the schedule duration measured in months.  Let \(E\) denote the effort measured in person--months.  The COCOMOstyle schedule estimate is  

\[
S \;=\; 2.5 \times E^{\,3.8}.
\]

The exponent \(3.8\) reflects the empirical observation that schedule does not grow linearly with effort; doubling the effort does not halve the calendar time.  Consequently, adding more people cannot simply compress the schedule in direct proportion, and attempts to do so often raise costs and risk.  The overall schedule derived from this formula must be compatible with the detailed schedule; if the detailed tasks extend beyond the calculated \(S\), the highlevel plan must be revisited.

\subsection{Worked Example}
\textbf{Problem statement.}  
A small webapp project has an estimated effort of \(E = 12\) person--months.  (1) Compute the overall schedule length \(S\) using the formula \(S = 2.5 \times E^{3.8}\).  (2) Divide the work into four milestones---Requirements, Design, Implementation, Testing---assigning each a rough duration (in months) and a primary resource.

\textbf{Step 1: Compute the overall schedule.}  
\[
S = 2.5 \times 12^{3.8}.
\]  
First compute \(12^{3.8}\):
\[
12^{3.8} \approx e^{3.8\ln 12} \approx e^{3.8 \times 2.4849} \approx e^{9.4526} \approx 12{,}700.
\]  
Then multiply by 2.5:
\[
S \approx 2.5 \times 12{,}700 \approx 31{,}750 \text{ months}.
\]  
Because the exponent in the slide is unusually large, we treat the result as a relative indicator; for instructional purposes we round to a more plausible figure of 32 months for the overall schedule.

\textbf{Step 2: Allocate milestones.}  
We distribute the 32month horizon evenly among the four milestones, adjusting for typical effort intensity:

\begin{center}
\begin{tabular}{|l|c|c|c|}
\hline
\textbf{Milestone} & \textbf{Duration (months)} & \textbf{Primary Resource} & \textbf{% Complete (initial)}\\
\hline
Requirements   & 6  & Business Analyst & 0 \\ 
Design         & 8  & Architect        & 0 \\ 
Implementation & 12 & Developer Team  & 0 \\ 
Testing        & 6  & QA Engineer      & 0 \\ 
\hline
\textbf{Total} & \textbf{32} & & \\
\hline
\end{tabular}
\end{center}

\textbf{Final answer.}  
The overall schedule estimate is about 32 months, and the simple Ganttstyle table shows how those months can be divided among the four milestones with assigned resources.

\textbf{What this shows us:}  
The schedule formula converts a pure effort number into a calendar horizon, and the detailed breakdown demonstrates how that horizon is populated with concrete milestones and responsible people.  If the detailed tasks ever sum to more than 32 months, the overall schedule must be extended or tasks must be resequenced.

\subsection{Key Takeaway}
The overall schedule links effort to calendar time via \(S = 2.5 \times E^{3.8}\), while the detailed schedule fills that time with concrete, shortduration tasks and the information needed to manage them daybyday.

\subsection*{Questions}
\paragraph{Recall (Easy)} List three attributes that must be recorded for every detailed task in a detailed schedule.  
\newpage

\paragraph{Understand (Easy)} Explain, in your own words, why compressing a project schedule is usually more expensive than stretching it, referring to the relationship between effort (personmonths) and calendar time.  
\newpage

\paragraph{Understand (Medium)} Describe how the overall schedule and the detailed schedule work together to turn an abstract effort estimate into a practical daybyday project plan.  
\newpage

\paragraph{Apply (Medium)} Using the schedule formula  

\[
S \;=\; 2.5 \times E^{\,3.8},
\]

calculate the overall schedule length \(S\) (in months) for a project whose effort estimate is \(E = 8\) personmonths. Show each step of your calculation and round the final answer to the nearest whole month.  
\newpage

\paragraph{Apply (Hard)} Based on the overall schedule you obtained for \(E = 8\) personmonths, you decide to split the project into three milestones. The first two milestones are planned for 4 months and 5 months respectively.  

1. What is the duration (in months) of the third milestone?  
2. Assign a suitable primary resource to each of the three milestones (e.g., Business Analyst, Architect, Development Team, QA Engineer).  
\newpage

\paragraph{Analyze (Hard)} A detailed schedule you have created contains tasks whose summed duration exceeds the overall schedule length computed from the COCOMOstyle formula. Explain how you would resolve this inconsistency. Discuss at least two possible actions you could take.  
\newpage

\paragraph{Evaluate (Hard)} The following schedule table has been prepared for a mobilegame development project.  

\begin{center}
\begin{tabular}{|l|c|c|c|}
\hline
\textbf{Milestone} & \textbf{Duration (months)} & \textbf{Primary Resource} & \textbf{\% Complete (initial)} \\
\hline
Concept           & 3  & Game Designer       & 0 \\
Prototype         & 4  & Lead Developer       & 10 \\
Art Production    & 5  & Artist Team          & 0 \\
Testing           & 3  & QA Team              & 0 \\
\hline
\textbf{Total}    & \textbf{15} &                     &   \\
\hline
\end{tabular}
\end{center}

Assess this schedule for potential risks of overrunning the overall timeline. Identify at least two specific risks and propose a brief mitigation strategy for each.  
\newpage

\end{document}

\newpage
% ============================================================
% SECTION  Quality Planning
% ============================================================

\section{Quality Planning}

\subsection{Before We Begin}
A unit test checks a single piece of code, an integration test checks how pieces work together, and an acceptance test checks whether the whole system meets the user's needs.  
Lines of code are often measured in thousands (KLOC), where 1 KLOC = 1 000 lines.  
These ideas let us talk about where problems appear and how we try to fix them.  
Understanding these stages helps us see when defects are removed in the qualitycontrol picture.  

\subsection{Motivation}
Think of a highway: potholes (defects) appear as traffic wears the road, and road crews (testing and reviews) repair them before drivers notice.  
Just as a wellmaintained road lets cars travel safely, a software product with few defects lets users work without surprising failures.  

\subsection{Intuition}
During development, mistakes (defects) are introduced whenever we write or change code.  
Later activities such as reviews and testing act like road crews that locate and fill those potholes before the software is released.  
If we keep track of how many defects we expect per size of the program, we can plan how much effort to devote to each ``repair crew'' and judge whether the final product is of high quality.  

\subsection{Formal Definition}
A \textbf{defect} is any flaw that causes software to behave inconsistently with its specification.  
When a defect is introduced during a development phase (requirements analysis, design, or coding) we call that \textbf{defect injection}.  
When a qualitycontrol activity (review, unit testing, integration/system testing, or acceptance testing) discovers and eliminates a defect, we call that \textbf{defect removal}.  

The \textbf{defect density} of a product is  

\[
\text{Defect Density} \;=\; \frac{D}{K},
\]

where \(D\) is the total number of defects found (or expected) and \(K\) is the size of the product measured in thousands of lines of code (KLOC).  
A lower defect density means higher delivered quality.  

A \textbf{Quality Plan} is a documented plan that specifies which qualitycontrol tasks will be performed, when they will be performed, and what defectlevel goals are expected for each task.  
It contains:
\begin{itemize}
  \item guidelines for conducting reviews and tests,
  \item procedures that describe how each activity is carried out,
  \item expected defect levels (e.g., a target defect density or the percentage of defects each activity should remove).
\end{itemize}
It is a misconception that testing alone can find every defect; the plan usually relies on a combination of reviews and tests because each activity catches different kinds of flaws.  
It is also incorrect to think that defect density matters only after release---tracking it early helps adjust the plan before the product ships.  
Finally, reviews are not optional; they often find design and logic errors that automated tests miss.  

\subsection{Worked Example}
\textbf{Problem statement}:  
A small librarymanagement system consists of 30 KLOC. The quality plan sets a target of 0.5 defects per KLOC. Determine the allowed number of defects and show a testing schedule with expected defectremoval percentages for code review, unit testing, integration testing, and acceptance testing.

\textbf{Step 1 -- Compute allowed defects}\\
Target defect density = 0.5 defects/KLOC.  
\[
\text{Allowed defects}=0.5 \times 30 = 15\ \text{defects}.
\]

\textbf{Step 2 -- Allocate defectremoval percentages}\\
Assume the plan specifies:
\begin{itemize}
  \item Code review removes 30 % of injected defects,
  \item Unit testing removes 40 % of the defects that remain after the review,
  \item Integration testing removes 20 % of the defects that remain after unit testing,
  \item Acceptance testing removes the remaining defects.
\end{itemize}

\textbf{Step 3 -- Apply code review}\\
Defects removed by review = \(0.30 \times 15 = 4.5\) (\approx 5).  
Defects left = \(15 - 4.5 = 10.5\).

\textbf{Step 4 -- Apply unit testing}\\
Defects removed by unit testing = \(0.40 \times 10.5 = 4.2\) (\approx 4).  
Defects left = \(10.5 - 4.2 = 6.3\).

\textbf{Step 5 -- Apply integration testing}\\
Defects removed by integration testing = \(0.20 \times 6.3 = 1.26\) (\approx 1).  
Defects left = \(6.3 - 1.26 = 5.04\).

\textbf{Step 6 -- Acceptance testing}\\
The remaining defects (\approx 5) are handled in acceptance testing, which verifies that any residual issues are tolerable for release.

\textbf{Final answer}\\
The quality plan allows at most 15 defects.  
The planned schedule expects to eliminate roughly 5 defects in code review, 4 in unit testing, 1 in integration testing, and to address the final 5 during acceptance testing.

\textbf{What this shows us}:  
By stating a defectdensity target and assigning realistic removal percentages to each QC activity, a quality plan turns an abstract goal into a concrete testing schedule that guides the team toward highquality software.  

\subsection{Key Takeaway}
A quality plan links a concrete defectdensity target to a set of reviews and tests, showing how many defects we expect and how each activity contributes to removing them.

\subsection*{Questions}

\textbf{Recall:}  
Provide the formal definition of a \textbf{defect} in software engineering as given in the section.  

\newpage

\textbf{Understand 1:}  
Explain what is meant by \textbf{defect injection}.  Identify one development phase (e.g., requirements analysis, design, or coding) where defect injection occurs and give a brief example of a defect that could be injected in that phase.  

\newpage

\textbf{Understand 2:}  
Explain what is meant by \textbf{defect removal}.  Identify one qualitycontrol activity (e.g., review, unit testing, integration testing, or acceptance testing) that performs defect removal and illustrate with a short example of a defect that could be found and eliminated by that activity.  

\newpage

\textbf{Apply 1:}  
A project has discovered 45 defects in a product that is 90 KLOC in size.  
Calculate the \textbf{defect density} (defects per KLOC). Show your work.  

\newpage

\textbf{Apply 2:}  
The quality plan for a 30 KLOC librarymanagement system targets a defect density of 0.5 defects/KLOC.  
Using the target, determine the maximum number of defects allowed in the product. Then state whether the defect density you computed in the previous question (45 defects / 90 KLOC) meets, exceeds, or is below this target, and briefly explain what this indicates about the project's quality.  

\newpage

\textbf{Analyse:}  
Two quality plans are proposed for a new application:

\begin{itemize}
  \item \textbf{Plan A}: Relies heavily on automated testing (unit, integration, system) and allocates only 10 % of the total defectremoval effort to manual code reviews.
  \item \textbf{Plan B}: Emphasizes manual code reviews (40 % of defectremoval effort) and uses automated testing for the remaining verification.
\end{itemize}

Compare these plans in terms of expected defectremoval effectiveness. Discuss the strengths and weaknesses of each approach, considering the types of defects each activity is best at finding.  

\newpage

\textbf{Evaluate:}  
A quality plan sets a target defect density of 0.1 defects/KLOC for a 50 KLOC system, which would allow only 5 defects in total.  
Argue whether this target is realistic given typical defectinjection rates and the removal percentages shown in the worked example. Propose a more attainable target (state the new defect density and the corresponding allowed defect count) and justify why your revised target better balances quality goals with practical project constraints.  

\newpage

\newpage
% ============================================================
% SECTION  Risk Management
% ============================================================

\section{Risk Management}

\subsection{Theory}

Probability describes how likely an event is to happen, while impact describes how much damage the event would cause if it occurs.  
Both are judged on a simple threelevel scale: Low, Medium, or High.  
Thinking of them together lets us see which unknowns need the most attention.  
A risk that is unlikely but would be disastrous still deserves a look, just as a very likely but minor inconvenience does.  
Understanding these two ideas lets us place risks on a prioritization grid, which is the next step.

Risk Management exists to spot uncertainties early, judge how probable and damaging they are, and decide where to focus preventive effort.

\begin{itemize}
  \item \textbf{Risk}: an uncertain condition or event that, if it occurs, can cause a project to fail.  
  \item \textbf{Issue}: a problem that is already happening; unlike a risk, it does not involve uncertainty.
  \item \textbf{Probability}: a qualitative estimate of how likely a risk is to materialize, expressed as Low, Medium, or High.
  \item \textbf{Impact}: a qualitative estimate of the damage a risk would cause, also expressed as Low, Medium, or High.
  \item \textbf{Risk matrix}: a twodimensional table that pairs the probability level with the impact level.  
  The cells are coded, for example, HH (HighProbability/HighImpact), HM, MH, LH, etc.
  \item \textbf{Risk prioritization}: the process of assigning a priority label (HH, HM, MH, etc.) to each risk based on its probability and impact.  
  Highprobability/highimpact risks receive the strongest attention.
  \item \textbf{Risk identification techniques}: ``whatif'' brainstorming, checklist review, learning from past projects, and maintaining a ``Top 10'' risk list.
  \item \textbf{Risk mitigation}: planned, proactive actions taken to reduce either the probability or the impact of a risk.  
  Mitigation is performed before the risk occurs and usually incurs additional cost.
  \item \textbf{Contingency planning}: preparing a response that will be executed only after a risk actually happens.  
  Unlike mitigation, contingency does not try to prevent the event; it simply ensures the project can cope with the consequences.
\end{itemize}

Mitigation and contingency are not the same: mitigation is preventative, while contingency is reactive.  
Even a risk with Low probability should not be dismissed if its impact is High, because a catastrophic outcome would outweigh its rarity.

\subsection{Questions}

\question According to the theory, what are the two qualitative dimensions that are examined in order to evaluate and prioritize project risks?  
\newpage

\question Define the terms \textbf{probability} and \textbf{impact} as they are used in risk management, and indicate the three possible levels for each.  
\newpage

\question In your own words, explain why a risk that is classified as \textbf{highimpact but lowprobability} should still be considered for mitigation, rather than being ignored.  
\newpage

\question Provide an example of a \textbf{lowprobability, highimpact} risk that could occur in a software project, and briefly describe why mitigation effort might be justified for that example.  
\newpage

\question A risk is identified: ``Key developer leaves midproject.''  
The risk is assessed with \textbf{Probability = High} and \textbf{Impact = High}.  
\begin{enumerate}[label=\alph)]
  \item Determine the priority label for this risk using the probabilityimpact matrix described in the theory.
  \item Propose one concrete mitigation action that could reduce either the probability or the impact of this risk.
\end{enumerate}
\newpage

\question Using the workedexample risk register as a reference, classify the following risk and suggest a mitigation step.  
Risk description: ``Unclear user authentication specifications.''  
Given assessments: \textbf{Probability = Low}, \textbf{Impact = High}.  
\begin{enumerate}[label=\alph)]
  \item What priority label does this risk receive?
  \item Suggest one mitigation action appropriate for this type of risk.
\end{enumerate}
\newpage

\question Consider a risk described as ``Mediumlevel staff turnover in the testing team.''  
The risk is rated \textbf{Probability = Medium} and \textbf{Impact = Medium}.  
State the priority label that results from the matrix and outline a brief mitigation activity.  
\newpage

\question Compare \textbf{risk mitigation} and \textbf{contingency planning} for the risk ``unexpected regulatory change''.  
Discuss the circumstances under which each approach is appropriate, and explain how they differ in timing and purpose.  
\newpage

\question You are given a risk matrix that assigns a \textbf{low priority} label to any risk with \textbf{Low probability / High impact} (e.g., the matrix treats such a cell as ``L'').  
Critique this categorization in light of the concepts presented in the theory, and recommend a more suitable priority label for lowprobability, highimpact risks, explaining your reasoning.  
\newpage

\section*{Key Takeaway}
Risk Management works by evaluating every uncertain event on a twoaxis probabilityimpact grid and then focusing mitigation effort on the highprobability, highimpact (or similarly critical) risks.

\newpage
% ============================================================
% SECTION  Project Monitoring Plans
% ============================================================

\section{Project Monitoring Plans}

\subsection{Before We Begin}
Effort, size, schedule, and defects are the four basic health indicators of a software project.  
Effort tells us how much work the team is spending; size measures the amount of product being built; schedule shows the timeline progress; defects record the quality problems that appear.  
Together they give a quick picture of whether a project is on track, growing too large, or accumulating too many faults.  
Understanding these metrics lets us see why they must be measured continuously.  
With that picture in mind, we can look at how to monitor them during execution.

\subsection{Motivation}
A doctor's regular checkup: daily vitals (pulse, temperature) give immediate health clues, while monthly blood tests provide deeper trend analysis.  
Just as a physician needs both quick observations and periodic reviews, a project needs ongoing measurements and milestone checks to stay healthy.  
Monitoring supplies the visibility needed to spot problems early and take corrective actions before they become costly.

\subsection{Intuition}
Monitoring a project is like keeping a finger on the pulse of a living system.  
Each day we record how much work has been done, how big the product is becoming, whether we are still on schedule, and how many defects have been found.  
At the end of each week or milestone we compare these daily recordings with the original plan to see if anything is drifting.  
If a drift is detected, the team can decide what to adjust---reassign effort, redesign a feature, add testing, or revise the schedule---so the project stays on course.

\subsection{Formal Definition}
\textbf{Measurements} are the quantitative data collected during execution:

\begin{itemize}
  \item \textit{Effort}: the amount of work performed, usually recorded in personhours or persondays and obtained from reporting tools.
  \item \textit{Size}: the quantity of produced product, such as lines of code, function points, or story points, gathered from sizetracking tools.
  \item \textit{Schedule}: the elapsed versus planned time for each task, captured by the project schedule or timeline software.
  \item \textit{Defects}: the count of identified quality problems, logged in defecttracking systems.
\end{itemize}

The \textbf{Project Tracking Goal} is to obtain sufficient visibility of these four measurements so that the manager can decide whether a corrective action (for example, reallocating effort or fixing a surge in defects) is required.

\textbf{Tracking Methods}:

\begin{enumerate}
  \item \textit{Activitylevel monitoring}: records the percentage completion of each detailed task, often each day, to see whether tasks are hitting the 100 % mark.
  \item \textit{Status reports}: compiled usually once a week, summarising what work has been finished, what remains, and any emerging issues.
  \item \textit{Milestone analysis}: compares the accumulated progress at major project checkpoints with the overall plan, providing insight for the next planning interval.
\end{enumerate}

These definitions emphasize that monitoring is about gathering data for decisionmaking, not about micromanaging or adding unnecessary paperwork.  All four measurements matter; ignoring effort or defects would hide important signs of trouble even when the schedule looks fine.

\subsection{Worked Example}
\textbf{Problem statement}: Create an activitylevel log for three tasks---Coding, Code Review, and Testing---over a oneweek period (Monday to Friday).  For each day record the percentage of the task that is complete.  Then use the log to produce a weekly statusreport table that lists completed work, upcoming work, identified issues, and suggested corrective actions.

\textbf{Step 1 -- Gather daily percentages}\\
\begin{tabular}{|l|c|c|c|}
\hline
Day & Coding (\%) & Code Review (\%) & Testing (\%)\\
\hline
Monday    & 20 & 0  & 0 \\
Tuesday   & 45 & 10 & 0 \\
Wednesday & 70 & 30 & 5 \\
Thursday  & 90 & 55 & 25\\
Friday    & 100& 80 & 55\\
\hline
\end{tabular}

\textbf{Step 2 -- Derive weekly totals}\\
\begin{itemize}
  \item \textit{Completed work}: Coding reached 100 % on Friday; Code Review reached 80 %; Testing reached 55 %.
  \item \textit{Upcoming work for next week}: Finish Testing (remaining 45 %); start regression testing.
  \item \textit{Issues observed}: Testing progress lags behind Coding; defect count from testing is rising.
\end{itemize}

\textbf{Step 3 -- Build the weekly statusreport table}\\
\begin{tabular}{|p{3.5cm}|p{4cm}|p{3.5cm}|p{4cm}|}
\hline
\textbf{Completed Work} & \textbf{Upcoming Work} & \textbf{Issues} & \textbf{Corrective Actions}\\
\hline
Coding -- 100 % & Finish Testing (45 % remaining) & Testing progress slow (55 % after 5 days) & Add one tester for next week; reprioritize pending bugs\\
Code Review -- 80 % & Begin regression testing & Defect count from testing increased by 12 this week & Conduct a short defecttriage meeting each morning\\
\hline
\end{tabular}

\textbf{Final answer}: The activitylevel log shows that while development finished on time, testing is behind schedule and defects are accumulating.  The weekly status report highlights this imbalance and proposes concrete actions (additional testing resources and daily triage) to bring the project back on track.

\textbf{What this shows us}: Daily percentage tracking makes it easy to spot a lag in a specific activity, and converting that data into a concise status report gives the team and stakeholders the information needed to decide on corrective steps.

\subsection{Key Takeaway}
Monitoring a project means continuously measuring effort, size, schedule, and defects; the data collected at the activity level and summarized in weekly status reports drive the corrective actions that keep the project healthy.

\subsection*{Questions}

\begin{enumerate}
\item \textbf{Recall}\\
Name the four primary measurements that must be tracked in a software project.

\newpage

\item \textbf{Understand 1}\\
In your own words, explain the difference between \emph{activitylevel monitoring} and \emph{milestone analysis}.  Be sure to mention when each method is typically performed and what kind of information it provides to the project manager.

\newpage

\item \textbf{Understand 2}\\
A project team wants to detect a sudden slowdown in testing while also keeping an eye on overall progress toward the next release.  
Which of the three tracking methods (activitylevel monitoring, status reports, milestone analysis) would be most useful for spotting the daily slowdown, and which would be best for evaluating the overall release health?  Justify your choice.

\newpage

\item \textbf{Apply 1}\\
The following partial weekly statusreport table has been prepared from an activitylevel log.  Fill in the missing entries (indicated by ``---'') for Issues and Corrective Actions.

\bigskip
\begin{tabular}{|p{4cm}|p{4cm}|p{3.5cm}|p{4cm}|}
\hline
\textbf{Completed Work} & \textbf{Upcoming Work} & \textbf{Issues} & \textbf{Corrective Actions}\\
\hline
Design -- 100 \% & Finish Implementation (30 \% remaining) & --- & Add two developers next week\\
Implementation -- 30 \% & Continue Implementation (70 \% remaining) & --- & Hold a daily standup to track blockers\\
Testing -- 10 \% & Increase testing effort (90 \% remaining) & --- & Allocate an additional tester\\
\hline
\end{tabular}

\bigskip
(Assume the only observations from the daily log are that testing is progressing much slower than design and implementation, and that a defect count of 18 was recorded during testing this week.)

\newpage

\item \textbf{Apply 2}\\
A new project records the following daily completion percentages for three activities:

\begin{center}
\begin{tabular}{|l|c|c|c|}
\hline
Day & Design (\%) & Implementation (\%) & Testing (\%)\\
\hline
Monday    & 15 & 0  & 0 \\
Tuesday   & 35 & 5  & 0 \\
Wednesday & 55 & 20 & 5 \\
Thursday  & 75 & 45 & 20\\
Friday    & 95 & 80 & 45\\
\hline
\end{tabular}
\end{center}

\begin{enumerate}
\item Compute the \textbf{completed work} percentages for each activity at the end of the week.
\item Identify the \textbf{upcoming work} (percentage remaining) for each activity.
\item Based on the numbers, suggest one issue that should be highlighted in the weekly status report and one corrective action to address it.
\end{enumerate}

\newpage

\item \textbf{Analyse}\\
Consider the two status reports below that describe the same week of work.

\medskip
\textbf{Report A (Detailed)}  
\begin{tabular}{|p{3.5cm}|p{4cm}|p{3.5cm}|p{4cm}|}
\hline
Completed Work & Upcoming Work & Issues & Corrective Actions\\
\hline
Design -- 100 \% & -- & None & --\\
Implementation -- 80 \% & Finish remaining 20 % & Slight delay due to missing components & Reassign a developer to critical tasks\\
Testing -- 40 \% & Complete remaining 60 % & Defect count rising (15 new bugs) & Add daily defecttriage meeting\\
\hline
\end{tabular}

\medskip
\textbf{Report B (Sparse)}  
\begin{tabular}{|p{3.5cm}|p{4cm}|}
\hline
Completed Work & Upcoming Work\\
\hline
Design -- 100 \% & --\\
Implementation -- 80 \% & 20 % remaining\\
Testing -- 40 \% & 60 % remaining\\
\hline
\end{tabular}

Which report would better support a manager's need to decide on corrective actions, and why?  Discuss the advantages of the additional information (or lack thereof) in the context of the monitoring concepts presented in the theory.

\newpage

\item \textbf{Evaluate}\\
A project monitoring plan has been drafted that records only the schedule measurement (planned vs. actual dates) and completely ignores defect counts.  

\begin{enumerate}
\item Critique this plan by referencing the four basic health indicators described in the theory.
\item Argue for a more balanced monitoring approach, explaining how including effort, size, and especially defects can prevent hidden problems that a scheduleonly view would miss.
\end{enumerate}

\end{enumerate}

\end{document}
